%----------------------------------------------------------------------------------------
% NOTATIE EN FIGUREN VAN DE CURSUS INTEGRALEN
% Staat in de preambule en niet in een hoofdstuk: ook de lijsten met oplossingen gebruiken dit.
%----------------------------------------------------------------------------------------
% Bewaarde grafiek, geplaatst door maak-pdf.py: bestand en pagina in figuren-cache/
\newcommand{\grafiekbestand}[2]{\includegraphics[page=#2,trim=3bp 3bp 3bp 3bp]{figuren-cache/#1}}

\newcommand{\dx}{\,\mathrm{d}x}
\newcommand{\dt}{\,\mathrm{d}t}
\newcommand{\du}{\,\mathrm{d}u}
\newcommand{\e}{\mathrm{e}}
\newcommand{\bgsin}{\operatorname{bgsin}}
\newcommand{\bgtan}{\operatorname{bgtan}}
\newcommand{\bgcos}{\operatorname{bgcos}}
% Rechter- en linkerlimiet, met de ongelijkheid onder de pijl: \limr{c}{0} en \liml{c}{1}
\newcommand{\limr}[2]{\lim\limits_{\substack{#1 \to #2 \\ >}}}
\newcommand{\liml}[2]{\lim\limits_{\substack{#1 \to #2 \\ <}}}
\newcommand{\ig}{\displaystyle\int}
\newcommand{\grens}[3]{\Bigl[#1\Bigr]_{#2}^{#3}}

% Zoals stappen, maar met een bredere kolom voor de berekening: voor lange formules
\newenvironment{stappenbreed}{%
  \par\renewcommand{\arraystretch}{1.25}%
  \noindent\begin{tabular}{@{}>{$\displaystyle}p{9.8cm}<{$}@{\hspace{4mm}}>{\RaggedRight\small\itshape\arraybackslash}p{4.9cm}@{}}%
}{\end{tabular}\par}

% Zoals stappen, met eigen kolombreedtes: \begin{stappenvrij}{6.2cm}{3.8cm}. Voor naast een figuur.
\newenvironment{stappenvrij}[2]{%
  \par\renewcommand{\arraystretch}{1.25}%
  \noindent\begin{tabular}{@{}>{$\displaystyle}p{#1}<{$}@{\hspace{3mm}}>{\RaggedRight\small\itshape\arraybackslash}p{#2}@{}}%
}{\end{tabular}\par}

% Tekst of berekening links, figuur rechts: \naast[breedte rechts]{links}{rechts}
% De breedte rechts komt uit de figuur zelf; het optionele argument wordt niet gebruikt.
\newsavebox{\naastbox}
\newcommand{\naast}[3][]{%
  \par\noindent\sbox{\naastbox}{#3}%
  \begin{minipage}[t]{\dimexpr\linewidth-\wd\naastbox-3mm\relax}\vspace{0pt}#2\end{minipage}\hfill
  \begin{minipage}[t]{\wd\naastbox}\vspace{0pt}\usebox{\naastbox}\end{minipage}\par\addvspace{1mm}}
% Een rij figuren, gecentreerd
\newcommand{\figuren}[1]{\par\addvspace{1mm}{\centering #1\par}\addvspace{1.5mm}}
% Opgave met een grafiek rechts: \begin{metgrafiek}[breedte]{tikz} opgave en deelvragen \end{metgrafiek}
% De solution komt na \end{metgrafiek}, zoals bij metprent.
\newsavebox{\grafiekbox}
\newenvironment{metgrafiek}[2][]{%
  \sbox{\grafiekbox}{#2}%
  \begin{minipage}[t]{\dimexpr\linewidth-\wd\grafiekbox-3mm\relax}%
}{%
  \end{minipage}\hfill
  \begin{minipage}[t]{\wd\grafiekbox}\vspace{-2mm}\usebox{\grafiekbox}\end{minipage}\par}

% Rooster met assen: \rooster{xmin}{xstap}{xmax}{ymin}{ystap}{ymax}, met getallen op de assen (gehele waarden).
% \roosterleeg doet hetzelfde zonder getallen, voor eigen labels.
\newcommand{\roosterleeg}[6]{%
  \draw[black!14,very thin,xstep=#2,ystep=#5] (#1,#4) grid (#3,#6);
  \draw[->,grijs] (#1,0) -- ({#3+.4*#2},0) node[right,text=black] {$x$};
  \draw[->,grijs] (0,#4) -- (0,{#6+.4*#5}) node[above,text=black] {$y$};}
\newcommand{\rooster}[6]{%
  \roosterleeg{#1}{#2}{#3}{#4}{#5}{#6}
  \pgfmathtruncatemacro{\rtweex}{#1+#2}\pgfmathtruncatemacro{\rtweey}{#4+#5}
  \foreach \t in {#1,\rtweex,...,#3}{\ifnum\t=0\else\node[below,inner sep=1.5pt,text=black] at (\t,0) {$\t$};\fi}
  \foreach \t in {#4,\rtweey,...,#6}{\ifnum\t=0\else\node[left,inner sep=1.5pt,text=black] at (0,\t) {$\t$};\fi}}

% In de functievoorschriften is% In de functievoorschriften is \x de veranderlijke; schrijf (\x)^2 en sin(\x r).
\newcommand{\assen}[4]{%
  \draw[->,lijngrijs] (#1,0) -- (#2,0) node[right,text=black] {$x$};
  \draw[->,lijngrijs] (0,#3) -- (0,#4) node[above,text=black] {$y$};}
\newcommand{\kromme}[4][]{\draw[accent,thick,#1] plot[domain=#2:#3,samples=90,smooth] (\x,{#4});}
% \gebied: tussen de grafiek en de horizontale as. \gebiedtussen: tussen twee grafieken.
\newcommand{\gebied}[4][accent!30]{\fill[#1] plot[domain=#2:#3,samples=70] (\x,{#4}) -- (#3,0) -- (#2,0) -- cycle;}
\newcommand{\gebiedtussen}[5][accent!30]{\fill[#1] plot[domain=#2:#3,samples=70] (\x,{#4}) -- plot[domain=#3:#2,samples=70] (\x,{#5}) -- cycle;}
% \riemann{aantal}{a}{b}{voorschrift}{plaats}: rechthoeken, met de hoogte op plaats 0 (links), 0.5 (midden) of 1 (rechts)
\newcommand{\riemann}[5]{%
  \pgfmathsetmacro{\rbreedte}{(#3-(#2))/#1}%
  \foreach \i in {1,...,#1}{%
    \pgfmathsetmacro{\xl}{#2+(\i-1)*\rbreedte}%
    \pgfmathsetmacro{\xr}{\xl+\rbreedte}%
    \pgfmathsetmacro{\x}{\xl+#5*\rbreedte}%
    \pgfmathsetmacro{\hoogte}{#4}%
    \draw[accent,fill=accent!22] (\xl,0) rectangle (\xr,\hoogte);}}
% \lichaam[kleur]{a}{b}{voorschrift}: het lichaam dat ontstaat door de grafiek om de x-as te wentelen, met de cirkel in b
\newcommand{\lichaam}[4][accent!30]{%
  \fill[#1] plot[domain=#2:#3,samples=50] (\x,{#4}) -- plot[domain=#3:#2,samples=50] (\x,{-(#4)}) -- cycle;
  \draw[accent,thick] plot[domain=#2:#3,samples=50,smooth] (\x,{#4});
  \draw[accent,thick] plot[domain=#2:#3,samples=50,smooth] (\x,{-(#4)});
  {\def\x{#3}\pgfmathsetmacro{\lstraal}{#4}\draw[accent,fill=accent!45] (#3,0) ellipse ({.14*\lstraal} and \lstraal);}}
